\scanpage{96}
en conjecturant que ces conditions \inserted{nécessaires} sont suffisamment draconiennes \deleted{pour} (même sans entrer ds la considération des $\Theta'$) pour assurer que $\Theta$ provient bien d'un $u\in\Gamma$, \deleted{\unsure{pour}} lorsque $X$ est (disons) une \uline{courbe algébrique anabélienne}, avec extension de type fini de $\mathbf{Q}$.
\margin{Je ne vois pas très bien si la condition a) \inserted{\unsure{fixe $x_0$ de point générique}}, sans b), devrait suffire pour assurer que $\Theta$ vient de $\Gamma$. Et puis, \unsure{$K$ est de t.f.} \unsure{sur tous ces objets, dans le groupe profini} \unsure{de $\Gamma$}.}
\margin{Il n'est pas évident que l'hypothèse anabélienne soit requise --- à examiner! (\unsure{N'est-ce pas?}). Pour $\mathbf{G}_m$, ou une courbe elliptique, la variation: $X/K$ \unsure{me semble également} \unsure{inessentielle}. \uline{En fait, c'est faux pour $\mathbf{G}_m$}.}
Bien sûr, la question de l'unicité de $u$ n'est pas \unsure{impliquée|indiquée} ici --- quand $K$ est une extension algébrique de $\mathbf{Q}$, on s'attend à l'\uline{unicité} de $u$. Quand $K$ est de degré de transcendance $>0$, alors pour pouvoir s'attendre à l'unicité de $u$, il faut au moins supposer que le corps modulaire de $X$ a ce degré de transcendance, i.e. que $K$ est une extension \uline{finie} de ce corps modulaire. Je ne serais pas étonné qu'on ait unicité dans ce cas.
\margin{Ça doit \unsure{alors se ramener au premier des cas, via l'application de} \unsure{Белый|Belyi}. \unsure{Pour l'unicité, bien sûr, la condition} \unsure{a)} \unsure{absolument analogue}\unsure{...!}}

Je reviens sur la question de la relevabilité d'un \deleted{$\unsure{\Theta}$} endomorphisme $\Theta$ d'un groupoïde connexe $\mathcal{T}$ à un revêtement $\widetilde{\mathcal{T}}$, quand $\Theta$ fixe un sommet
\mspage{95}
